\documentclass[11pt,a4paper]{article}
\usepackage[T1]{fontenc}
\usepackage{lmodern,microtype,amsmath,amssymb,amsthm,booktabs}
\usepackage[margin=25mm]{geometry}
\usepackage[colorlinks=true,linkcolor=blue,urlcolor=blue]{hyperref}
\newtheorem{theorem}{Theorem}
\newtheorem{lemma}[theorem]{Lemma}
\newtheorem{corollary}[theorem]{Corollary}
\newcommand{\E}{\mathbb E}
\newcommand{\cc}{C^{\mathrm{causal}}}
\title{A Bounded Causal Entropy Tax\\\large Unequal block lengths remove the growing boundary cost}
\author{Mathematical construction for the C3-C model}
\date{26 September 2026}
\begin{document}
\maketitle
\begin{abstract}
For any two finite output laws with strictly different Shannon entropies,
we construct exact finite-horizon causal realizations whose seed entropy is
$n\max\{H(P),H(Q)\}+O(1)$. The constant is independent of the horizon and
the construction is valid for every adaptive action policy. An explicit
example has a strictly positive residual bounded by $14$ bits for every
$n$. The construction uses geometrically increasing and decreasing block
lengths in the anti-diagonal arrangement. This result concerns a strict
entropy gap; it does not resolve the equal-entropy logarithmic problem.
\end{abstract}

\section{Statement and model}
All logarithms and entropies are base two. A discrete seed fixes a
history-indexed deterministic response tree. At each step an action selects
one of two finite laws $P,Q$. Under every adaptive policy the outputs must
have exactly the prescribed conditional laws. Let $\cc_n(P,Q)$ be the
infimum of the seed entropy. This is the C3-C model of the supplied
\emph{Who Pays the Bill?}, Version 4; no synchronous cross-history identity
is assumed. Zero-probability symbols may be discarded.

\begin{theorem}[Bounded tax under a strict entropy gap]\label{thm:general}
If $H(P)>H(Q)$, there is a finite constant $K(P,Q)$ such that, for all $n$,
\[
\boxed{nH(P)\le\cc_n(P,Q)\le nH(P)+K(P,Q).}
\]
Moreover, $\cc_n-nH(P)$ is nondecreasing and has a finite limit.
The same assertion holds after exchanging $P,Q$.
\end{theorem}

\begin{theorem}[An explicit nonzero bounded residual]\label{thm:explicit}
Let $P$ be uniform on eight symbols and let $Q=(1/3,2/3)$. For every
positive integer $n$,
\[
\boxed{\frac{h_2(1/3)}8\le\cc_n(P,Q)-3n<14,}
\qquad \frac{h_2(1/3)}8=0.1147869792568\ldots,
\]
where $h_2$ is binary entropy. Thus the residual has order $\Theta(1)$,
with a strictly positive finite limiting value.
\end{theorem}

The exact value of that limit is not determined here. The upper constant
$14$ is deliberately loose. Different action alphabets cause no difficulty;
one may pad $Q$ with zeros if a common alphabet is required.

\section{The construction: variable-length anti-diagonal blocks}
Choose positive integers $s_1,\ldots,s_m$ with sum $n$. Partition both
length-$n$ streams $X,W$ into consecutive blocks of these lengths. Draw
the following components independently:
\[
X^{(1)}\sim P^{\otimes s_1},\qquad
W^{(m)}\sim Q^{\otimes s_m},\qquad
(W^{(j)},X^{(j+1)})\sim J_j\quad(1\le j<m),
\]
where $J_j$ is any coupling of
$Q^{\otimes s_j}$ and $P^{\otimes s_{j+1}}$.
On the $a$th occurrence of action $P$, output $X_a$; on the $b$th
occurrence of action $Q$, output $W_{n+1-b}$.

\begin{lemma}[Causal validity and entropy accounting]\label{lem:anti}
This construction is exact under every adaptive policy. If
$H(J_j)\le s_{j+1}H(P)+d_j$, then its seed satisfies
\begin{equation}\label{eq:account}
H(S)\le nH(P)+s_mH(Q)+\sum_{j=1}^{m-1}d_j.
\end{equation}
\end{lemma}
\begin{proof}
Put $T_j=s_1+\cdots+s_j$. Seeing any coordinate of $X^{(j+1)}$
requires at least $T_j+1$ occurrences of $P$. Seeing any coordinate of
$W^{(j)}$ in reverse order requires at least $n-T_j+1$ occurrences of
$Q$. Seeing both therefore requires $n+2$ actions, which is impossible.
For a fixed action word, each independent seed component contributes
coordinates from at most one of its marginal product blocks. The observed
outputs consequently have the required product law. An output branch of
an adaptive deterministic policy fixes one such action word; applying the
same statement to that branch proves exactness. Randomized policies follow
by conditioning on their independent randomness.

Independence of the seed components gives
$H(S)=s_1H(P)+s_mH(Q)+\sum_j H(J_j)$, proving \eqref{eq:account}.
\end{proof}

Unlike equal-length blocking, this formula only pays for the \emph{last}
$Q$-block as an uncovered boundary. Both boundary blocks can remain small
while the interior block lengths become large. A strict entropy gap permits
the necessary decrease of block lengths on the far side of the construction.

\section{A complete elementary proof of the explicit example}
\begin{lemma}[Uniform fine grid]\label{lem:grid}
If $A$ is uniform on $M$ symbols and $B$ has $K\ge2$ symbols, there is
an exact coupling with
\[
H(A,B)\le\log_2 M+\frac{K-1}{M}\log_2 K.
\]
\end{lemma}
\begin{proof}
Use the common-uniform interval coupling on $[0,1)$. The $A$-intervals
all have length $1/M$. Only cells containing an internal endpoint of a
$B$-interval can have positive conditional entropy $H(B\mid A=a)$.
There are at most $K-1$ such cells, and each contributes at most
$(1/M)\log_2 K$. Endpoints on grid boundaries do not split cells.
\end{proof}

Apply the lemma to $A=P^{\otimes p}$, $B=Q^{\otimes q}$. If $p\ge q/2$,
\begin{equation}\label{eq:explicit-block}
H(A,B)-3p\le q\,2^{q-3p}\le f(q),
\qquad f(q)=q\,2^{-q/2}.
\end{equation}
No approximation of either marginal is involved.

\paragraph{An integer mountain for every horizon.}
For $n\ge2$, choose the unique $k\ge0$ such that
$4\cdot2^k\le n+2<8\cdot2^k$, and put $s=2^k$.
Use the left tail $(1,2,\ldots,2^{k-1})$ and its reverse as the right
tail; for $k=0$ the tails are empty. The remaining length is
\[
R=n-2(s-1),\qquad 2s\le R<6s.
\]
Set $t=\lceil R/s\rceil$, and divide $R$ as evenly as possible into $t$
integer central blocks. There are at most six central blocks. Each has
length between $s/2$ and $s$ (and length $1$ if $s=1$).
The resulting sequence sums to $n$, begins and ends at $1$, and satisfies
\[
s_{j+1}\ge s_j/2\quad\hbox{for every adjacent pair.}
\]
Thus \eqref{eq:explicit-block} applies to every coupling in
Lemma~\ref{lem:anti}. Each dyadic tail length appears at most twice, so
\begin{equation}\label{eq:14}
\cc_n-3n\le H(Q)+2\sum_{j=0}^\infty f(2^j)
                        +6\max_{q\in\mathbb N}f(q)<14.
\end{equation}
Here are elementary certificates for the last inequality. The sequence
$f(q)$ is maximized at $q=3$, since
$f(q+1)/f(q)=(1+1/q)/\sqrt2$. Hence
\[
\max_q f(q)=\frac{3}{2\sqrt2}<\frac{16}{15}.
\]
The first five terms of the dyadic series are
$1/\sqrt2,1,1,1/2,1/16$. The remaining tail is smaller than $1/1024$
(its first term is $1/2048$, with subsequent ratios below $1/2$). Therefore
\[
\sum_{j\ge0}f(2^j)
<\frac{71}{100}+1+1+\frac12+\frac1{16}+\frac1{1024}
<\frac{33}{10}.
\]
Since $H(Q)<1$, the right side of \eqref{eq:14} is strictly below
$1+2(33/10)+6(16/15)=14$. The unrounded expression in
\eqref{eq:14} is approximately $13.82244702$.

\paragraph{Strict positivity, not a relabeling trick.}
At horizon one, a seed can be replaced by its pair of action outputs
without increasing entropy. Write $x_i$ for the joint probability of
the $i$th $P$-symbol and the first $Q$-symbol. Then
\[
0\le x_i\le1/8,\qquad \sum_{i=1}^8x_i=1/3,\qquad
H(P,Q)=3+\frac18\sum_{i=1}^8h_2(8x_i).
\]
A concave function on this polytope has a minimizer at an extreme point.
There, at most one $x_i$ is strictly between $0$ and $1/8$.
Since $8/3=2+2/3$, the unique fractional row has conditional probability
$2/3$. Consequently
\[
\cc_1=3+\frac18h_2(1/3).
\]
For any realization of horizon $n+1$, first choose action $P$. Conditional
on that first output, the remaining seed is an admissible horizon-$n$
seed. The chain rule gives
\begin{equation}\label{eq:monotone}
\cc_{n+1}\ge H(P)+\cc_n.
\end{equation}
This proves the lower bound in Theorem~\ref{thm:explicit}, and also handles
$n=1$ in the upper bound. A positive, nondecreasing, bounded residual has
a positive finite limit. This completes the explicit example.

\section{Why every strict entropy gap suffices}
Write $h_P=H(P)>h_Q=H(Q)>0$. Choose a constant
$1<\Lambda<h_P/h_Q$. The case $h_Q=0$ is simpler: $Q$ is deterministic,
and a length-$n$ iid $P$ stream gives exactly $\cc_n=nh_P$.

\begin{lemma}[Exponentially small block excess]\label{lem:exp}
There are constants $C,c>0$, depending only on $P,Q,\Lambda$, such that
for all positive integers $p,q$ with $1/\Lambda\le p/q\le\Lambda$,
an exact coupling satisfies
\[
H(P^{\otimes p},Q^{\otimes q})\le ph_P+F(q),
\qquad F(q)=C(q+1)e^{-cq}.
\]
\end{lemma}
\begin{proof}
Choose $\varepsilon>0$ with
$(h_P-\varepsilon)/\Lambda>h_Q+\varepsilon$.
Except on an event exponentially small in $p$, a $P$-word has
probability at most $u=2^{-p(h_P-\varepsilon)}$.
Except on an event exponentially small in $q$, a $Q$-word belongs to a
set of at most $K=2^{q(h_Q+\varepsilon)}$ words.
These exponential estimates follow directly from the exponential Markov
inequality for sums of bounded independent information contents: the
centered log moment-generating function has derivative zero at the origin,
so any fixed nonzero deviation has a strictly positive exponential rate.
If an information content is constant, its exceptional event is empty.

Arrange the typical $Q$-words consecutively in the common-uniform interval
coupling. Decode a $Q$-word from a $P$-word whenever the latter's interval
is contained in a typical $Q$-interval; otherwise choose an arbitrary
decoder output. Decoder error can occur only on an atypical $P$-word,
an atypical $Q$-word, or a typical $P$-interval crossing one of the at
most $K+1$ typical-region boundaries. Thus its probability is at most
\[
\delta\le e^{-c_1p}+e^{-c_2q}+(K+1)u\le C_1e^{-c_3q},
\]
with constants adjusted if an exceptional event is empty.
For the actual error probability $e\le\delta$, the elementary decoding
entropy bound gives
\[
H(Q^{\otimes q}\mid P^{\otimes p})
\le h_2(e)+e\,q\log_2|\operatorname{supp}Q|.
\]
This is bounded by $C(q+1)e^{-cq}$ for all sufficiently large $q$.
Increase $C$ to cover the finitely many smaller integer pairs. The coupling
itself has exact marginals throughout; typical sets are used only to
estimate its entropy.
\end{proof}

\paragraph{A mountain with arbitrarily gentle slopes.}
Choose $1<\lambda<\Lambda$. Pick positive integers $M,b$ so large that
\[
\lambda+1/b\le\Lambda,\qquad
\frac{M}{M+1}-\frac1b\ge\frac1\Lambda,
\]
and put $a_j=\lceil b\lambda^j\rceil$. Then
$a_{j+1}/a_j\le\Lambda$ and $a_j\ge b\lambda^j$.
For $n\ge Mb$, choose the largest integer $k\ge0$ with
\[
n\ge2\sum_{j<k}a_j+Ma_k.
\]
Use the two tails $(a_0,\ldots,a_{k-1})$ and its reverse.
For the remaining length $R$, maximality gives
\[
Ma_k\le R<2a_k+Ma_{k+1}\le(2+M\Lambda)a_k.
\]
Divide $R$ evenly into $t=\lceil R/a_k\rceil$ integer central blocks.
Each central length $\ell$ satisfies
\[
\frac{a_k}{\Lambda}\le
a_k\frac{M}{M+1}-1\le\ell\le a_k.
\]
There are at most $T=\lceil2+M\Lambda\rceil$ such blocks.
All successive ratios of the entire sequence lie in
$[1/\Lambda,\Lambda]$. Its last block is $b$ if $k\ge1$, and at most
$b$ if $k=0$.

Apply Lemmas~\ref{lem:anti} and~\ref{lem:exp}. Uniformly in $n\ge Mb$,
\begin{equation}\label{eq:generalK}
\cc_n-nh_P\le
bh_Q+2\sum_{j=0}^\infty F(a_j)+T\sup_{q\ge1}F(q)<\infty.
\end{equation}
The series converges because $a_j$ grows geometrically and $F$ decreases
exponentially at infinity. For the finitely many horizons $n<Mb$, use
independent iid streams and enlarge the constant by at most $Mb h_Q$.
The constant-action $P$ path supplies the lower bound $nh_P$.
Equation~\eqref{eq:monotone} gives monotonicity, completing the proof of
Theorem~\ref{thm:general}.

\section{What this does and does not settle}
The strict entropy gap supplies enough excess information inside a long
$P$-block to absorb a somewhat longer $Q$-block at exponentially small
additional entropy cost. Growing and then shrinking the blocks leaves
bounded endpoint lengths and a summable sequence of coupling costs.
This is the step that equal-size blocks miss.

When $H(P)=H(Q)$, the interval $(1,H(P)/H(Q))$ is empty and the argument
cannot be used. In particular, it does not override the square-root
lower bound for the equal-entropy, unequal-varentropy crash-test pair,
nor prove a logarithmic bound for a variance-matched equal-entropy pair.
The constant in \eqref{eq:generalK} need not remain bounded as the entropy
gap tends to zero. The constructed seed may depend on $n$; no single
projectively consistent infinite seed is asserted.

The anti-diagonal mechanism is inherited from Section 18 of the
user-supplied manuscript by J. Mikul\'ik, \emph{Who Pays the Bill?},
Version 4 (19 August 2026). This note supplies the variable-length
construction and its entropy estimates with complete proofs. No claim
of historical priority is made. The accompanying Python script verifies
integer mountain partitions, exact rational block marginals, and the
rational certificates for the explicit bound; the all-horizon and adaptive
claims are proved above rather than inferred from those tests.
\end{document}
